%% Flint Secure — NCIT Initial Project Proposal (Guidelines 2019)
\documentclass[proposal,minor]{ncitproject}

\ProjectTitle{Flint Secure\\[0.35em]{\large Real-Time Fraud Detection and Device Intelligence Platform}}
\DegreeName{Bachelor of Engineering in Computer Engineering}
\DepartmentName{Computer Engineering}
\SupervisorName{[Supervisor Name]}
\SupervisorDesignation{[Designation]}
\SubmissionDate{04 Jun 2026}

\SetStudents{%
  Aaditya Yadav & 231201 \\
  Apil Khadka & 231210 \\
  Ujashna Dangol & 231247 \\
}

\begin{document}

\maketitlepage

\begin{ncitabstract}
Nepal's digital payment sector is expanding quickly, but most payment service providers still
depend on static, rule-based fraud checks. Flint Secure is proposed as a real-time fraud
detection and device intelligence platform calibrated for Nepal's market, where SIM swap
attacks rose 70\% year-on-year (FIU-Nepal, 2024) and FATF compliance pressure requires
stronger controls.

The platform will expose \texttt{/v1/identify} for stable device identity across software
updates (three-layer fingerprint matching) and \texttt{/v1/score} for composite risk over
seven dimensions: device trust, amount, time, velocity, network, identity, and recipient.
Each transaction will receive a score from 0--100 and a decision of ALLOW, FLAG, or BLOCK,
with a design target of under 50\,ms P95 latency at roughly 3,200 requests per second.

Implementation is planned in Go with dual Redis instances, NATS JetStream, PostgreSQL,
multi-platform SDKs, and Python workers for profiling and trust updates. Success will be
measured through load testing and a 168-case automated test suite with at least 92\% pass
rate, aligned with Pokhara University minor-project requirements and NCIT proposal-defense
criteria.
\end{ncitabstract}

\Keywords{fraud detection, device fingerprinting, real-time risk scoring, fintech security, Nepal digital payments, FATF compliance, SIM swap detection, behavioral profiling}

%% Reserved — not included in the NCIT proposal PDF (guideline has no ``document type'' notice).

\section{Problem Statement}

Over the past five years, Nepal's digital payment ecosystem has grown rapidly. Mobile wallets
such as eSewa, Khalti, and IME Digital now serve millions of users, and the Nepal Rastra Bank
reports more than ten million registered mobile-banking users (2024). This growth expands
the attack surface for financial fraud while most payment service providers (PSPs) still rely
on simple rule engines---thresholds on amount or frequency---that cannot detect coordinated
attacks.

Three gaps define the problem for this project:
\begin{enumerate}
  \item \textbf{No adaptive detection.} Rule sets must be changed manually and cannot learn
        from new fraud patterns such as SIM swap, credential stuffing, or device emulation.
  \item \textbf{No device intelligence.} Transactions are not linked to a stable physical
        device identity, so one-device--many-account (money mule) behaviour is invisible.
  \item \textbf{No behavioural baseline.} Systems do not maintain per-user profiles for
        amount, time-of-day, or recipient patterns, so anomalies are missed until losses occur.
\end{enumerate}

Nepal's FATF grey-list placement and FIU-Nepal's report of a 70\% year-on-year rise in SIM
swap fraud (2024) show that existing controls are insufficient. The result is high false
negatives (fraud undetected) and, when rules are tightened blindly, high false positives that
block legitimate customers and reduce PSP revenue.

\section{Project Objectives}

The primary objective is to \textbf{design, implement, and validate} Flint Secure as a
production-oriented fraud platform for Nepali PSPs. The following specific objectives will
guide the semester work:
\begin{enumerate}
  \item Develop a \textbf{device identity layer} using multi-signal fingerprinting and
        drift-tolerant fuzzy matching so devices remain recognisable after browser or OS updates.
  \item Build a \textbf{composite risk engine} with seven weighted signal dimensions calibrated
        using Nepal-specific fraud indicators (including SIM swap prevalence).
  \item Meet a \textbf{performance target} of P95 latency below 50\,ms and sustained throughput
        of at least 3,200 scoring requests per second on commodity cloud hardware.
  \item Implement \textbf{SIM swap detection} through hashed phone-to-device binding with
        first-write-wins semantics and hard-floor blocking rules.
  \item Deliver an \textbf{asynchronous analytics pipeline} (NATS + Python workers) for user
        profiling, device trust, and recipient risk without blocking the scoring path.
  \item Prove correctness via a \textbf{168-test automated suite} with pass rate of at least 92\%.
\end{enumerate}

\section{Significance of the Study}

Commercial platforms (Stripe Radar, Sift, Sardine) are built for Western transaction volumes
and pricing models that are out of reach for typical Nepali PSPs. They also lack calibration
for domestic vectors such as SIM swap-heavy account takeover. Flint Secure is significant
because it proposes \textbf{affordable, locally relevant} fraud infrastructure using open
components (Go, Redis, NATS, PostgreSQL) deployable on Oracle Cloud Free Tier ARM64.

For the college and university, the project will aim to show that minor-project outcomes can
meet industry-style non-functional targets (latency, throughput, security middleware) while
remaining reproducible by a small student team. For the national context, it will support
digital financial inclusion and reduces reliance on foreign fraud services that are neither
priced nor tuned for Nepal.

\section{Literature Study/Review}

\subsection{Device Fingerprinting}
Eckersley (2010) showed that browser fingerprints are highly unique in large samples, forming
the empirical basis for passive device identification. Laperdrix et al.\ (2016) demonstrated
that canvas rendering traces GPU hardware and informs our Tier-1 canvas hash. Vastel et al.\
(2018) tracked fingerprint evolution over time; their work supports our proposed three-layer
cascade: exact hash match, stored-ID recovery after drift, and fuzzy similarity with
evolution penalties.

\subsection{Real-Time Fraud Detection and Scoring}
West and Bhattacharya (2016) review machine-learning approaches to payment fraud; ensemble
methods consistently outperform single models. Fiore et al.\ (2019) address severe class
imbalance via synthetic fraud samples---relevant for early deployment when labels are scarce.
Dal Pozzolo et al.\ (2015) provide a framework for probability calibration that informs our
planned ALLOW/FLAG/BLOCK thresholds. For the initial release we propose transparent rule-based
scoring; ML integration is deferred to future work (see \texttt{archive/future-work.tex}).

\subsection{Velocity Analysis and SIM Swap}
Bolton and Hand (2002) identify velocity features (frequency, amount variance, recipient
diversity) as among the strongest real-time predictors. Our design maps these to Redis sorted-set
counters queried in one pipelined round-trip. Mjaaland et al.\ (2021) document SIM swap in
mobile banking; FIU-Nepal (2024) confirms rapid growth in Nepal, justifying a high identity-risk
weight in the proposed scoring model.

\subsection{Comparison with Existing Products}
Stripe Radar, Sift, and Sardine offer mature cloud scoring but at per-event costs unsuitable
for many domestic PSPs and without Nepal-specific weighting. Flint Secure is positioned as
\textbf{open, self-hostable infrastructure} with explicit calibration for local fraud data.

\section{Proposed Methodology}

\subsection{Overview and Architecture}
Flint Secure will follow a layered architecture: HTTP middleware and handlers (request path),
pure business logic (scoring engine), and asynchronous persistence (NATS workers to
PostgreSQL). The scoring path will prefetch all signals from Redis, call a pure
\texttt{scoring.Score()} function with no I/O, return the decision, then publish events and
velocity updates in background goroutines. This separation is proposed to satisfy the
sub-50\,ms latency objective.

\subsection{Middleware and Security}
Each request will pass through: request-ID assignment, client IP extraction, GeoIP enrichment
(MaxMind GeoLite2), API-key authentication (Redis cache with PostgreSQL fallback), tiered rate
limiting (per IP, key, and device), and optional HMAC-SHA256 signing with Redis-backed replay
protection ($\pm$5 minute timestamp window).

\subsection{Storage Design}
\begin{itemize}
  \item \textbf{Redis Persist} --- device records, user profiles, phone bindings, monthly quotas
        (AOF enabled; 90-day device TTL).
  \item \textbf{Redis Cache} --- velocity sorted sets, rate-limit counters, API-key cache (no
        persistence; optimised for throughput).
  \item \textbf{PostgreSQL} --- transactions, alerts, clients, API-key hashes (writes only via
        workers, not on the hot path).
  \item \textbf{NATS JetStream} --- subjects \texttt{device.first\_seen} and
        \texttt{transactions.scored} linking Go producers to Python consumers.
\end{itemize}

\subsection{Device Identification (\texttt{/v1/identify})}
The SDK will hash collected signals client-side (\texttt{local\_fp}). The server will apply:
\textbf{Layer~1} exact Redis lookup (expected $\sim$85\% of calls); \textbf{Layer~2} recovery via
stored device ID when fingerprints drift; \textbf{Layer~3} fuzzy match over 14 tiered signals
within a block-key bucket (screen size and OS family). New devices will be enrolled and indexed
for later matches.

\subsection{Risk Scoring (\texttt{/v1/score})}
Seven components---device, amount, time, velocity, network, identity, recipient---will be
weighted and combined with interaction terms (e.g.\ new device $\times$ large amount).
Hard-floor rules will force BLOCK for definitive signals (SIM swap mismatch, datacenter IP,
impossible device consistency). Phone numbers will be stored only as SHA-256 hashes with
first-write-wins device binding.

\subsection{Async Worker Pipeline}
Python workers will consume NATS events to: index devices and transactions in PostgreSQL;
build user amount percentiles and active-hour profiles; update device trust scores; aggregate
recipient risk from prior BLOCK decisions; refresh client configuration caches in Redis.

\subsection{Technology Stack and SDKs}
Go (chi router), dual Redis, NATS, PostgreSQL, Python~3 workers; JavaScript, Android, and iOS
SDKs for signal collection; k6 for load testing; deployment target Oracle Cloud ARM64 (4 OCPU
instance in the validation plan).

%% Design diagrams (non-floating) inside Methodology
\subsection{System Design Diagrams}

The following figures summarise the proposed architecture.

\newpage
\umlfigure[1.9]{uml-use-case}{use\_case}{Updated use case diagram (PSP Backend, Web SDK, Fraud Analyst).}
Figure~\ref{fig:uml-use-case} outlines the primary use cases of the system. The Web SDK collects client-side device parameters and sends them for identification. The PSP Backend initiates transaction risk scoring to decide payment actions. The Fraud Analyst monitors triggered alerts, performs manual reviews, and assigns fraud labels to train risk models.
\newpage

\umlfigure[1.9]{uml-seq-identify}{sequencediag}{Sequence diagram 1: device identification flow.}
Figure~\ref{fig:uml-seq-identify} illustrates the sequence of operations for device identification. The client's Web SDK extracts browser parameters, computes a local fingerprint hash, and submits them to the API Gateway. The Gateway executes a multi-step matching routine (exact, cached ID lookup, and fuzzy heuristics) to resolve or generate a unique Device ID, returning it to the SDK.
\newpage

\umlfigure[1.9]{uml-seq-score}{sequencediag2}{Sequence diagram 2: scoring flow with Redis and NATS JetStream.}
Figure~\ref{fig:uml-seq-score} details the real-time scoring sequence. When a PSP requests a score, the API Gateway queries the fast Redis cache in a single batch to retrieve risk limits, velocity counters, and device history. The gateway runs the rule-based scoring engine to determine the action (ALLOW, FLAG, or BLOCK), returns the decision immediately, and asynchronously publishes a transaction event to NATS JetStream.
\newpage

\begin{landscape}
\umlfigure[1.6]{uml-seq-three}{sequencediag3}{Sequence diagram 3: additional sequence flow.}
Figure~\ref{fig:uml-seq-three} depicts the asynchronous processing flow handled by background workers. Background workers subscribe to NATS JetStream topics to process scored transactions. They update transaction metrics in PostgreSQL, update device profiling counters in Redis, and flag anomalies without blocking the critical transaction execution path.
\end{landscape}
\newpage

\umlfigure[1.9]{uml-seq-four}{sequencediag4}{Sequence diagram 4: additional sequence flow.}
Figure~\ref{fig:uml-seq-four} illustrates the workflow for alert review and risk model updates. High-risk transactions trigger alerts that are queued for review. Analysts inspect these alerts and apply manual fraud labels. These labels are used to retrain and evaluate machine learning models. Once deployed, the updated models dispatch webhooks to the registered merchant endpoints.
\newpage

\umlfigure[1.0]{er-diagram}{er_diagram}{Entity-Relationship (ER) Diagram of the database schema.}
Figure~\ref{fig:er-diagram} shows the Entity-Relationship schema for the system. It represents the structural relations between core tables: Merchants register Webhooks and own Devices. Devices identify Transactions, which in turn are scored by Risk Scores. Transactions also trigger Alerts reviewed by Analysts, who generate Fraud Labels to train the Machine Learning Models that dispatch Webhooks.



\section{Proposed Performance Analysis Methodology and Validation Scheme}

Validation will be \textbf{planned and repeatable}, not anecdotal. Table~\ref{tab:perf} lists
metrics, methods, and acceptance thresholds tied to Objectives~3 and~6.

\begin{table}[ht]
  \centering
  \caption{Proposed performance and quality validation}
  \label{tab:perf}
  \small
  \begin{tabularx}{\textwidth}{@{}lXcc@{}}
    \toprule
    \textbf{Metric} & \textbf{Procedure} & \textbf{Target} & \textbf{Tool} \\
    \midrule
    P95 \texttt{/score} latency & 10-min sustained load, co-located Redis & $<$50\,ms & k6 \\
    Sustained throughput & Ramp to peak RPS, watch error rate & $\geq$3,200 req/s & k6 \\
    P95 \texttt{/identify} (L1/L3) & Isolated layer tests & $<$500\,ms & Go tests \\
    HTTP 5xx under load & Monitor server logs & $<$0.1\% & k6 + logs \\
    Functional coverage & 168 automated cases & $\geq$92\% pass & Go/Python \\
    SIM swap floor & Inject phone rebinding & BLOCK triggered & Unit tests \\
    \bottomrule
  \end{tabularx}
\end{table}

\textbf{Weight calibration} will use three inputs: FIU-Nepal fraud-vector statistics (SIM swap
share), Nepal Rastra Bank payment-distribution data (amount and hour), and 2,400 labelled
pilot transactions (Q3--Q4 2024) supplied by a domestic PSP partner. Latency profiling will
separate Redis pipeline time from scoring CPU time to confirm the pure-function design.

\textbf{Defense readiness} will include a 10-minute presentation, a supervisor meeting log
(per college Appendix~B), and a live or recorded demo showing \texttt{ALLOW}, \texttt{FLAG},
and \texttt{BLOCK} decisions on sample transactions.

\section{Proposed Deliverable/Output}

\subsection{At proposal defense (this document)}
\begin{itemize}
  \item This initial proposal report (8--10 pages, NCIT format).
  \item UML~2.0 models: use case, component, and sequence diagrams (PlantUML source in repository).
  \item Ten-minute proposal presentation and supervisor meeting log (Appendix~B, college guideline).
\end{itemize}

\subsection{Planned by end of semester}
\begin{table}[ht]
  \centering
  \caption{Planned outputs after proposal approval (not yet built)}
  \label{tab:deliverables}
  \small
  \begin{tabularx}{\textwidth}{@{}lXl@{}}
    \toprule
    \textbf{No.} & \textbf{Output} & \textbf{Target} \\
    \midrule
    D1 & Go REST API (\texttt{/v1/identify}, \texttt{/v1/score}) & Week~8 \\
    D2 & JavaScript, Android, iOS SDKs & Week~9 \\
    D3 & Redis schemas + PostgreSQL migrations & Week~7 \\
    D4 & NATS worker pipeline (6 consumers) & Week~10 \\
    D5 & 168-test suite + k6 load scripts & Week~12 \\
    D6 & API \& deployment documentation & Week~13 \\
    D7 & Mid-term and final reports + demo & Defenses \\
    \bottomrule
  \end{tabularx}
\end{table}

\section{Project Task and Time Schedule}

This schedule begins at \textbf{proposal defense} (start of semester). Mid-term and final
defenses follow NCIT timing: at least four weeks after proposal, and at least four weeks
after mid-term, respectively.

\begin{table}[ht]
  \centering
  \caption{Project task and time schedule (14-week semester)}
  \label{tab:schedule}
  \small
  \begin{tabularx}{\textwidth}{@{}clXl@{}}
    \toprule
    \textbf{Wk} & \textbf{Milestone} & \textbf{Tasks} & \textbf{Output} \\
    \midrule
    1 & \textbf{Proposal defense} & Submit this proposal, UML, presentation; supervisor log & This report \\
    2--3 & Foundation & Repo setup, middleware, Redis/Postgres, CI & D3 \\
    4--5 & Identify API & Fingerprint cascade, block keys, SDK alpha & D1, D2 \\
    6--7 & Score API & Velocity, scoring engine, SIM swap, quota & D1 \\
    8 & \textbf{Mid-term defense} & Design review, progress demo, supervisor log & Mid-term report \\
    9--10 & Async layer & NATS subjects, Python workers, profiling & D4 \\
    11--12 & Validation & k6 runs, 168 tests, bug fixes, tuning & D5 \\
    13--14 & \textbf{Final defense} & Final report, demo, documentation, presentation & D6, D7 \\
    \bottomrule
  \end{tabularx}
\end{table}

% Keep floats (UML figures) above references — do not let them drift past bibliography
\clearpage
\section{Bibliography/References}

\begin{thebibliography}{99}

\bibitem{bolton2002}
R.~J. Bolton and D.~J. Hand,
\newblock Statistical fraud detection: A review,
\newblock \emph{Statistical Science}, 17(3):235--255, 2002.

\bibitem{dalpozzolo2015}
A.~Dal Pozzolo, O.~Caelen, R.~A. Johnson, and G.~Bontempi,
\newblock Calibrating probability with undersampling for unbalanced classification,
\newblock \emph{2015 IEEE Symposium Series on Computational Intelligence}, 2015.

\bibitem{eckersley2010}
P.~Eckersley,
\newblock How unique is your web browser?
\newblock \emph{Privacy Enhancing Technologies Symposium (PETS 2010)}, EFF, 2010.

\bibitem{fiore2019}
U.~Fiore, A.~De Santis, F.~Perla, P.~Zanetti, and F.~Palmieri,
\newblock GAN-based augmentation for credit-card fraud detection,
\newblock \emph{Information Sciences}, 479:448--455, 2019.

\bibitem{fiu2024}
Financial Intelligence Unit Nepal,
\newblock \emph{Annual Report 2023/24},
\newblock Nepal Rastra Bank, Kathmandu, 2024.

\bibitem{laperdrix2016}
P.~Laperdrix, W.~Rudametkin, and B.~Baudry,
\newblock Beauty and the beast: Diverting modern web browsers to build unique fingerprints,
\newblock \emph{IEEE Symposium on Security and Privacy}, 2016.

\bibitem{mjaaland2021}
B.~B. Mjaaland, J.~Heeger, and S.~D. Wolthusen,
\newblock SIM swap: A new insider attack in the mobile banking ecosystem,
\newblock \emph{IEEE CSR}, 2021.

\bibitem{vastel2018}
A.~Vastel, P.~Laperdrix, W.~Rudametkin, and R.~Rouvoy,
\newblock FP-STALKER: Tracking browser fingerprint evolutions,
\newblock \emph{IEEE Symposium on Security and Privacy}, 2018.

\bibitem{west2016}
J.~West and M.~Bhattacharya,
\newblock Intelligent financial fraud detection: A comprehensive review,
\newblock \emph{Computers \& Security}, 57:47--66, 2016.

\end{thebibliography}


\end{document}
